\documentclass[12pt]{article}
\usepackage[brazil]{babel}
\usepackage[utf8]{inputenc}
\usepackage[T1]{fontenc}
\usepackage{amsmath, amssymb}
\usepackage{geometry}
\usepackage{graphicx}
\usepackage{xcolor}

\definecolor{maincolor}{RGB}{0,100,50}
\usepackage{fancyhdr}
\definecolor{forestgreen}{RGB}{0,100,50}
\usepackage[most]{tcolorbox}
\newtcolorbox{block}[1]{%
  colback=green!5,
  colframe=maincolor,
  coltitle=white,
  title=#1,
  fonttitle=\bfseries,
  boxrule=1.5pt,
  arc=3mm,
  left=6pt,
  right=6pt,
  top=6pt,
  bottom=6pt
}

\geometry{margin=2.5cm}
\pagestyle{fancy}
\fancyhf{}

\renewcommand{\headrulewidth}{2pt}
\renewcommand{\footrulewidth}{2pt}

\renewcommand{\headrule}{\hbox to\headwidth{%
  \color{forestgreen}\leaders\hrule height \headrulewidth\hfill}}

\renewcommand{\footrule}{\hbox to\headwidth{%
  \color{forestgreen}\leaders\hrule height \footrulewidth\hfill}}

\setlength{\headheight}{32pt}

\fancyhead[L]{%
  \includegraphics[height=1.2cm]{Universidade_de_Rio_Verde.jpg}
}

\fancyhead[R]{%
  \raggedleft
  \textcolor{forestgreen}{%
    \small
    \textbf{Lista de Exercícios — Divisibilidade}\\
    Me.~Sergio Souza Novak\\
    \textit{Lógica para Engenharia de Software}
  }
}

\fancyfoot[C]{%
  \textcolor{forestgreen}{\thepage}
}

\begin{document}

\vspace*{0.2cm}

\section*{Instruções}

\begin{block}{Utilize a definição formal de divisibilidade}
Sejam $a,b \in \mathbb{Z}$, com $b \neq 0$. Dizemos que $b$ divide $a$ se existe
$k \in \mathbb{Z}$ tal que
\[
a = bk.
\]
Neste caso, escrevemos $b \mid a$.
\end{block}

Justifique rigorosamente todas as afirmações utilizando apenas definições formais
e propriedades dos números inteiros.
% =======================
% EXERCÍCIOS (SEM RESPOSTA)
% =======================

\section*{Exercícios}

\begin{enumerate}

\item Demonstre que, se $a \mid b$ e $a \mid c$, então $a \mid (b+c)$.

\item Demonstre que, se $a \mid b$, então $a \mid (-b)$.

\item Prove que, se $a \mid b$ e $b \mid c$, então $a \mid c$.



\end{enumerate}

\newpage

% =======================
% GABARITO
% =======================

\section*{Gabarito — Respostas}
\begin{block}{1. Soma preserva divisibilidade}
Suponha que $a \mid b$ e $a \mid c$. Pela definição de divisibilidade, existem
inteiros $k,m \in \mathbb{Z}$ tais que $b = ak$ e $c = am$. Somando essas
igualdades, obtemos $b + c = ak + am = a(k+m)$
Como $k+m \in \mathbb{Z}$, segue diretamente da definição que
$a \mid (b+c)$, concluindo a demonstração.
\end{block}

\begin{block}{2. Divisibilidade pelo oposto}
Suponha que $a \mid b$. Então, pela definição de divisibilidade, existe
$k \in \mathbb{Z}$ tal que $b = ak$. Multiplicando ambos os lados da igualdade
por $-1$, temos
$-b = a(-k).$
Como $-k \in \mathbb{Z}$, conclui-se que $a \mid (-b)$, conforme a definição
formal de divisibilidade.
\end{block}

\begin{block}{3. Transitividade da divisibilidade}
Suponha que $a \mid b$ e $b \mid c$. Pela definição, existem inteiros
$k,m \in \mathbb{Z}$ tais que $b = ak$ e $c = bm$. Substituindo a expressão
de $b$ na igualdade que define $c$, obtemos
$
c = (ak)m = a(km).
$
Como $km \in \mathbb{Z}$, segue pela definição de divisibilidade que
$a \mid c$, o que demonstra a transitividade da relação de divisibilidade.
\end{block}
\end{document}

